\begin{figure}[htbp]
\centering
\begin{tikzpicture}[alt={Two actual light-theme Bitmap Display observations of a sixteen by eight storage view. A double outline marks the logical eight by six frame in its upper-left region; the gray right half and bottom two rows are padding. First panel shows placement minus two comma one, with eight opaque copies among twelve visible cells. Second panel shows placement one comma one, with thirteen opaque copies among twenty visible cells. Native marker one identifies the copied word 01FF00FF in both panels; marker two identifies the opaque zero word in the second panel.}]
\node[anchor=north west,inner sep=0,font=\small\bfseries] (atitle) at (0,0) {1. Placement $(-2,1)$: left columns clipped};
\node[anchor=north west,inner sep=0] (a) at ([yshift=-2mm]atitle.south west) {\MARSRevisionCrop{revision/captures/clipped/clipped-left.png}{0.82\linewidth}{0.032031250}{0.155555556}{0.032031250}{0.178888889}};
\node[anchor=north west,inner sep=0,font=\small\bfseries] (btitle) at ([yshift=-4mm]a.south west) {2. Placement $(1,1)$: complete source visible};
\node[anchor=north west,inner sep=0] (b) at ([yshift=-2mm]btitle.south west) {\MARSRevisionCrop{revision/captures/clipped/clipped-full.png}{0.82\linewidth}{0.032031250}{0.155555556}{0.032031250}{0.178888889}};
\foreach \panel in {a,b} {
\begin{scope}[shift={(\panel.south west)},x={($(\panel.south east)-(\panel.south west)$)},y={($(\panel.north west)-(\panel.south west)$)}]
\draw[black,line width=1.8pt] (0,.25) rectangle (.5,1);
\draw[white,line width=.7pt] (0,.25) rectangle (.5,1);
\end{scope}}
\foreach \panel/\x/\y/\n in {a/.09375/.6875/1,b/.28125/.6875/1,b/.15625/.6875/2} {
\begin{scope}[shift={(\panel.south west)},x={($(\panel.south east)-(\panel.south west)$)},y={($(\panel.north west)-(\panel.south west)$)}]
\node[font=\scriptsize\bfseries,fill=white,text=black,draw=black,circle,inner sep=1pt] at (\x,\y) {\n};
\end{scope}}
\end{tikzpicture}
\caption{Actual completed copies. The native publication outline separates logical frame cells from allocation padding. Marker 1 identifies an opaque word whose low RGB bytes match the key; marker 2 identifies an opaque zero that draws black. Key cells retain the previous background. All 128 native display cells were checked for each observation.}
\label{fig:revision-clipped-observations}
\end{figure}
